% CHECKPOINT36: formalization of frozen actual_elliptic_original_cusp_bridge35.md
% SHA256 ec8fef231a71e2329ed2cfbdedaafe90055dce8a45a2ef77aa428eb1a85eadde.
\statement{A.122}{Proved}{The exact elliptic boundary and original third-cusp bridge}
Let \(\alpha\) be the original negative involution of full
\(\widehat{\mathrm{GT}}\), with original parameter \(\phi\), canonical
commutator word \(u\), and actual \(d=1+2k\in\widehat{\mathbb Z}^{\times}\)
with the same even \(k\) as above. In \(E=\widehat F(X,Y)\), put
\[
 p=[Y,X],\qquad q=X,\qquad z=(pq)^{-1}=YX^{-1}Y^{-1},\qquad
 H=\overline{\langle p,q\rangle}.
                                                                    \tag{A.122.1}
\]
The pair \((p,q)\) is a free profinite basis of the actual embedded
subgroup \(H\), with its full profinite topology.
Write
\[
 \begin{gathered}
 u_0=u(p,q),\quad u_2=u(z,p),\quad U=u(q,p)=u_0^{-1},\quad V=u_2,\\
 \Phi=\phi(q,p),\qquad Q=\phi(z,p),\\
 \xi=U X^dU^{-1},\qquad
 \omega=p^kVYU^{-1},\qquad \eta(X)=\xi,\quad\eta(Y)=\omega .
 \end{gathered}                                                     \tag{A.122.2}
\]
Let \(\Theta\) be the actual elliptic involution supplied by A.121
and the full profinite section, and let
\(c_{\mathrm{ell}}(X)=X^{-1},\
c_{\mathrm{ell}}(Y)=p^{-1}Y\).
Then
\[
 \Theta(\xi)=\xi^{-1},\qquad \Theta(\omega)=p^{-d}\omega,       \tag{A.122.3}
\]
and the following three conditions are equivalent:
\[
 \begin{gathered}
 \Theta\eta=\eta c_{\mathrm{ell}},\qquad
 [\omega,\xi]=p^d,\\
 \psi_d(z)=\operatorname{Ad}\bigl(p^ku(z,p)\bigr)(z^d),
 \quad
 \psi_d(p)=p^d,\quad\psi_d(q)=u_0^{-1}q^du_0 .
 \end{gathered}                                                     \tag{A.122.4}
\]
The third condition is a specified conjugacy at the original third
cusp of \(H\). The closed normal subgroup of \(E\) generated by
\([\omega,\xi]p^{-d}\) is \(\Theta\)-stable. None of the equivalent
conditions is asserted to hold.

\begin{proof}
\emph{1. The actual embedded free pair.}
The open subgroup \(V_2\) defined by the \(Y\)-exponent modulo two
has the explicit Schreier basis
\[
 X,\quad YXY^{-1},\quad Y^2 .
\]
Thus its first two basis elements generate a free factor. Replacing
\(YXY^{-1}\) by \((YXY^{-1})X^{-1}=p\), while retaining \(q=X\),
is an invertible Nielsen basis change. Consequently
\(H=\widehat F(p,q)\) is an actual free factor of \(V_2\).
Its intrinsic full topology is its induced topology, by the
free-factor retraction. No completion of an infinite-index kernel
has been identified.

\emph{2. Restrict the actual involution and the original norm.}
At \(\lambda=-1\), the full profinite section gives
\[
 \Theta(X)=\Phi X^{-1}\Phi^{-1},\qquad
 \Theta(Y)=p^{-1}QY\Phi^{-1},\qquad
 \Theta(p)=p^{-1}.                                           \tag{A.122.5}
\]
The final identity is the boundary identity in the full proof of
\citet[Proposition~2.20]{A_Enriquez2014}.
This uses the actual full-GT pair \((-1,\phi)\); no section is
being applied to the candidate pair \((d,u)\).

On \((p,q)\), (A.122.5) is the original negative action with
parameter \(\phi(p,q)=\Phi^{-1}\). Apply the original norm
\(u=\phi\alpha(u)\) to \(u(p,q)=U^{-1}\) and invert. This gives
\[
 \Theta(U)=U\Phi^{-1}.                                       \tag{A.122.6}
\]
Direct substitution also gives
\(\Theta(z)=p^{-1}Qz^{-1}Q^{-1}p\).
Hence \(\operatorname{Ad}(p)\Theta\) restricts on \((p,z)\)
to the original negative action with parameter
\(\phi(p,z)=Q^{-1}\). Applying the same norm and duality to
\(u(p,z)=V^{-1}\) proves
\[
 \Theta(V)=p^{-1}VQ^{-1}p.                                   \tag{A.122.7}
\]
These are evaluations of the original canonical word; no
replacement norm solution or transporter has been chosen.

Substitution in the stated order gives the first equality of
(A.122.3), and for the second gives
\[
 \begin{aligned}
 \Theta(p^kVYU^{-1})
 &=p^{-k-1}VQ^{-1}p\,p^{-1}QY\Phi^{-1}\Phi U^{-1}\\
 &=p^{-k-1}VYU^{-1}=p^{-d}\omega .
 \end{aligned}
\]
The map \(c_{\mathrm{ell}}\) is the actual section of \((-1,1)\)
and is an involution; its second generator image is also
\(XYX^{-1}\).
Since \(\eta(p)=[\omega,\xi]\), the two sides of
\(\Theta\eta=\eta c_{\mathrm{ell}}\) agree on \(X\), and agree
on \(Y\) precisely when
\[
 p^{-d}\omega=[\omega,\xi]^{-1}\omega.
\]
This proves its equivalence with the boundary equation, without
cancelling through \(\eta\) or assuming it surjective.

\emph{3. The original third cusp, with its complete conjugator.}
The equality \(YqY^{-1}=pq=z^{-1}\) gives
\[
 \begin{aligned}
 [\omega,\xi]
 &=p^ku_2Y u_0(u_0^{-1}q^du_0)u_0^{-1}Y^{-1}
                         u_2^{-1}p^{-k}\xi^{-1}\\
 &=p^ku_2z^{-d}u_2^{-1}p^{-k}\xi^{-1}.
 \end{aligned}                                               \tag{A.122.8}
\]
All powers are actual continuous profinite powers. Therefore
\[
 \begin{aligned}
 [\omega,\xi]=p^d
 &\Longleftrightarrow
 p^ku_2z^{-d}u_2^{-1}p^{-k}=p^d\xi=\psi_d(z^{-1})\\
 &\Longleftrightarrow
 \psi_d(z)=\operatorname{Ad}(p^ku_2)(z^d).
 \end{aligned}
\]
This proves the last equivalence in (A.122.4) in the embedded
original free pair, with its specified conjugator retained.

\emph{4. Stability of the exact obstruction quotient.}
Put \(D=[\omega,\xi]p^{-d}\) and let \(J\) be its closed normal
closure in \(E\). Modulo \(J\), (A.122.3) gives
\[
 \begin{aligned}
 \Theta(D)
 &=[p^{-d}\omega,\xi^{-1}]p^d\\
 &=[[\omega,\xi]^{-1}\omega,\xi^{-1}]p^d
   =[\omega,\xi]^{-1}p^d=1 .
 \end{aligned}
\]
The middle simplification is a literal group identity. Thus
\(\Theta(J)\subseteq J\); involutivity of \(\Theta\) gives
equality. This constructs the obstruction quotient and does
not prove \(D=1\).

The pentagon-only cubic of A.121 supplies no omitted step in this
boundary equation. The section has been used for \(\alpha\)
only. In particular neither an elliptic action for the GT-zero
lift of A.118, the original full hexagon, nor a coherent-frame
or CM value is inferred.
\end{proof}
